\section{Robot Pack Schema Summary}
\label{app:schema}

This appendix summarizes the implemented format. The normative contract remains
\pathcode{docs/robot_pack_spec.md} and the Pydantic models under
\pathcode{src/modsim/robot_packs/schema.py}.

\begin{longtable}{L{0.24\textwidth}L{0.25\textwidth}L{0.39\textwidth}}
\caption{Robot Pack 0.1 principal documents and fields.}
\label{tab:schema-summary}\\
\toprule
Entity & Principal fields & Purpose \\
\midrule
\endfirsthead
\toprule
Entity & Principal fields & Purpose \\
\midrule
\endhead
Manifest & schema version, ID/name/version, metadata, views, assets, split paths,
mappings & Root identity and cross-document catalog \\
Asset manifest & URDF list, mesh directories, MuJoCo assets, Isaac assets, thumbnails &
Declared mechanical/visual resources \\
Module type & asset reference, root link, mass, joints, connectors, capabilities &
Reusable semantic hardware unit \\
Joint & source name, type, parent/child, axis, control modes, limits &
Semantic articulation and controller contract \\
Connector instance & type, parent link, frame/local pose, docking/approach axes,
metadata & Physical port placement \\
Connector type & activity, gender, compatibility, orientations, acceptance, physical
intent, limits, policy, metadata & Shared mating and lifecycle contract \\
Capability & kind, description, connector requirements & Advertised action/behavior \\
Model-view recipe & builder, modes, default, configuration & Portable view-generation request \\
Backend mapping & asset catalog and module link/joint/frame/constraint maps &
Semantic-to-backend name translation \\
\bottomrule
\end{longtable}

\subsection{Identifier grammar}

Most semantic IDs match
\code{\^{}[a-z][a-z0-9]*(?:\_[a-z0-9]+)*\$} and have a maximum length of 64.
Runtime compound identifiers use \code{<module>/<local>} for connectors/joints and a
canonical sorted pair for connection IDs. Display names are not IDs.

\subsection{Units}

Fields include units in their names wherever ambiguity is likely:
\code{mass\_kg}, \code{xyz\_m}, \code{rpy\_rad},
\code{position\_tolerance\_m}, \code{orientation\_tolerance\_rad},
\code{max\_relative\_velocity\_m\_s}, \code{effort\_nm}, and related forms.
Runtime commands and snapshots use SI.

\subsection{Connector policy defaults}

The effective default is explicit command required, measured alignment, no redock
cooldown, and no break force. Defaults are serialized and round-trip tested so Studio
editing does not silently erase connector-type metadata or policy.

\section{Runtime Events and State Transitions}
\label{app:events}

\begin{longtable}{L{0.27\textwidth}L{0.58\textwidth}}
\caption{Canonical semantic event taxonomy.}
\label{tab:event-taxonomy}\\
\toprule
Event & Meaning \\
\midrule
\endfirsthead
\toprule
Event & Meaning \\
\midrule
\endhead
\code{DockCandidateDetected} & A compatible pair entered candidate evaluation; emitted
on transition rather than every step. \\
\code{DockCommitted} & Backend created a physical constraint and the logical connection
was committed. \\
\code{DockFailed} & An explicit attempt failed compatibility, acceptance, guard, or
backend execution. \\
\code{AssemblyMerged} & A committed dock joined two previously separate components. \\
\code{UndockCommitted} & Backend removed the constraint and logical connection. \\
\code{UndockFailed} & Requested release did not execute. \\
\code{AssemblySplit} & An undock separated one connected component into multiple components. \\
\code{ConnectorOverloaded} & Reported physical load exceeded the weaker break-force policy. \\
\bottomrule
\end{longtable}

The actively used connector lifecycle states include free, candidate detected,
in-acceptance, latching, docked, releasing, and failed. The broader vocabulary also
contains aligning and load-bearing states for future richer execution. A failed
transient is cleared on the next backend ingest, while cooldown prevents immediate
retries where configured.

\subsection{Successful dock pseudocode}

\begin{lstlisting}[language=Python]
snapshot = adapter.snapshot()
world.ingest(snapshot)

for pair in sorted(candidate_pairs(world)):
    compatibility = check_compatibility(pair)
    acceptance = evaluate_acceptance(pair, world)
    guards = evaluate_guards(pair, world, commands)
    if not (compatibility.ok and acceptance.ok and guards.ok):
        report_commanded_failure_if_needed()
        continue

    request = build_connection_request(pair, acceptance)
    outcome = adapter.create_physical_connection(request)
    if outcome.accepted:
        world.apply(DockCommitted(..., handle=outcome.handle))
        world.apply(AssemblyMerged(...))  # only if components differed
    else:
        world.apply(DockFailed(reason="backend_refused", ...))
\end{lstlisting}

\section{SMORES-EP Parameters and Provenance}
\label{app:parameters}

\begin{longtable}{L{0.33\textwidth}L{0.22\textwidth}L{0.31\textwidth}}
\caption{Key \smoresep{} model and controller values.}
\label{tab:parameter-provenance}\\
\toprule
Quantity & Value & Provenance/status \\
\midrule
\endfirsthead
\toprule
Quantity & Value & Provenance/status \\
\midrule
\endhead
Module mass & \SI{0.4386368}{kg} & Fusion-derived geometry/inertial export \\
Wheel radius & \SI{0.040}{m} & Fusion-derived \\
Track width & \SI{0.0672}{m} & Fusion-derived \\
Wheel angular-speed cap & $\pi/2\,\mathrm{rad/s}$ & Published platform value \\
Wheel effort cap & \SI{0.04}{N.m} & Provisional simulation value \\
PAN/TILT effort cap & \SI{0.10}{N.m} & Provisional simulation value \\
Wheel velocity gain & $0.1\,\mathrm{N,m/(rad/s)}$ & Provisional simulation value \\
Hold $P,D$ & $1.0,0.02$ & Provisional simulation value \\
Solver/controller period & \SI{0.002}{s} default & Provisional tuned value \\
Maximum accepted timestep & \SI{0.005}{s} & Scenario safety bound \\
Connector position tolerance & \SI{0.006}{m} & Conservative semantic proxy \\
Pair near-contact gate & \SI{0.001}{m} & Two-module simulation choice \\
Physical-snake latch gate & \SI{0.006}{m} & Uses pack acceptance envelope \\
Orientation tolerance & \SI{25}{\degree} & Literature-motivated, not fully calibrated \\
Relative-speed tolerance & \SI{0.05}{m/s} & Pack semantic value \\
Redock cooldown & \SI{0.1}{s} & Pack semantic value \\
Release delay & \SI{0.08}{s} & Literature-derived switching delay model \\
Normal/shear/bending limits & 100 N / 50 N / 2 N m & Explicitly unverified/provisional \\
Tire friction tuple & $(0.05,1,0.01,10^{-4},10^{-4})$ & Provisional anisotropic contact \\
Skid friction tuple & $(0.08,0.08,0.005,10^{-4},10^{-4})$ & Provisional contact \\
Rigid-allocation residual cap & \SI{0.08}{m/s} & Provisional controller safety bound \\
\bottomrule
\end{longtable}

\section{Driver-to-Snake Topology and Routes}
\label{app:routes}

\subsection{Initial topology}

The six initial connector edges are:

\begin{enumerate}
  \item \code{module\_1/bottom}--\code{module\_2/pan};
  \item \code{module\_2/bottom}--\code{module\_3/pan};
  \item \code{module\_2/right}--\code{module\_4/left};
  \item \code{module\_4/right}--\code{module\_5/left};
  \item \code{module\_5/pan}--\code{module\_6/bottom}; and
  \item \code{module\_5/bottom}--\code{module\_7/pan}.
\end{enumerate}

\subsection{Physical route parameters}

Waypoints are $(f,l,h,d,\epsilon)$: forward offset, left offset, heading offset,
drive direction, and position tolerance, all relative to the action's final moving-root
pose. Heading tolerance is $\pi$ in the current route helper unless explicit alignment
is enabled.

\begin{longtable}{L{0.08\textwidth}L{0.72\textwidth}L{0.10\textwidth}}
\caption{Authored target-relative Driver-to-Snake routes.}
\label{tab:route-details}\\
\toprule
Action & Waypoints $(f,l,h,d,\epsilon)$ & Approach \\
\midrule
\endfirsthead
\toprule
Action & Waypoints $(f,l,h,d,\epsilon)$ & Approach \\
\midrule
\endhead
1 & $(0.34,0,0,+1,0.03)$; $(0.34,0.13,\pi/2,+1,0.03)$;
$(-0.06,0.13,\pi,+1,0.03)$; $(-0.06,0,-\pi/2,+1,0.03)$;
$(-0.03,0,0,+1,0.01)$ & forward \\
2 & $(-0.34,0,0,-1,0.03)$; $(-0.34,-0.13,-\pi/2,+1,0.03)$;
$(0.06,-0.13,0,+1,0.03)$; $(0.06,0,\pi/2,+1,0.03)$;
$(0.12,0,0,+1,0.015)$ & reverse \\
3 & $(-0.27,0.088,0,-1,0.03)$; $(-0.03,0,-0.35,+1,0.01)$ & forward \\
4 & $(0.27,-0.088,0,+1,0.075)$; $(0.03,0.008,-0.35,-1,0.01)$ & reverse \\
\bottomrule
\end{longtable}

The action-4 \SI{8}{mm} prebias and \SI{75}{mm} clearance tolerance are empirical
simulation tuning. They should be replaced by planned, feedback-corrected paths in a
general controller.

\section{Reproduction Commands}
\label{app:commands}

\subsection{Environment}

\begin{lstlisting}[language=bash]
uv sync --locked --extra dev --extra studio --extra mujoco --python 3.12
uv run --no-sync modsim --version
\end{lstlisting}

Equivalent editable pip installation:

\begin{lstlisting}[language=bash]
python3.12 -m venv .venv
.venv/bin/python -m pip install -e ".[studio,mujoco,dev]"
\end{lstlisting}

\subsection{Pack inspection and Studio}

\begin{lstlisting}[language=bash]
modsim pack inspect examples/robot_packs/smores_ep
modsim pack validate examples/robot_packs/smores_ep
modsim pack validate examples/robot_packs/smores_ep --profile simulation
modsim studio examples/robot_packs/smores_ep
\end{lstlisting}

\subsection{Physical two-module demonstration}

\begin{lstlisting}[language=bash]
modsim runtime examples/robot_packs/smores_ep \
  --backend mujoco \
  --demo smores_diff_drive_dock_undock \
  --model-view smores_topology
\end{lstlisting}

\subsection{Kinematic Driver-to-Snake}

\begin{lstlisting}[language=bash]
modsim runtime examples/robot_packs/smores_ep \
  --backend mujoco \
  --demo smores_driver_to_snake \
  --model-view smores_topology \
  --duration 14 \
  --dt 0.002 \
  --no-gravity
\end{lstlisting}

\subsection{Physical Driver-to-Snake}

\begin{lstlisting}[language=bash]
modsim runtime examples/robot_packs/smores_ep \
  --backend mujoco \
  --demo smores_physical_driver_to_snake \
  --model-view smores_topology \
  --speed 4
\end{lstlisting}

Omit \code{--speed} for real-time pacing. The factor is fixed at launch. Both windows
retain the final state until Stop or close.

\subsection{Verification}

\begin{lstlisting}[language=bash]
uv lock --check
uv run --no-sync ruff check .
uv run --no-sync ruff format --check .
uv run --no-sync pyright
uv run --no-sync pyright --project pyright-mujoco.json
uv run --no-sync pyright --project pyright-studio.json
uv run --no-sync pytest
\end{lstlisting}

Focused physical tests:

\begin{lstlisting}[language=bash]
uv run --no-sync pytest tests/test_mujoco_physical_docking.py
uv run --no-sync pytest tests/test_mujoco_physical_reconfiguration.py
uv run --no-sync pytest tests/test_mujoco_contact_mechanics.py
\end{lstlisting}

\section{Requirement-to-Evidence Traceability}
\label{app:traceability}

\begin{longtable}{L{0.11\textwidth}L{0.35\textwidth}L{0.40\textwidth}}
\caption{Traceability from \cref{tab:requirements} to principal implementation evidence.}
\label{tab:traceability}\\
\toprule
Req. & Implementation & Principal regression areas \\
\midrule
\endfirsthead
\toprule
Req. & Implementation & Principal regression areas \\
\midrule
\endhead
R1--R2 & schema, loader, validator & schema, loader, validator, SMORES example pack \\
R3 & backend base/registry, mock, MuJoCo adapter & registry, backend conformance \\
R4 & broad phase, acceptance, guards, docking & acceptance and docking suites \\
R5 & connection request/outcome and event commit & docking, MuJoCo weld, refusal tests \\
R6 & world revisions, assemblies, model topology & assemblies, revisions, model views \\
R7 & joint IDs/commands/session validation, MuJoCo effort & joint commands, MuJoCo backend \\
R8 & inspection DTO/protocol/presenter/process/window & runtime inspection/process/widgets/worker \\
R9 & safe paths, writer staging/rollback, Studio project & loader, writer, Studio project \\
R10 & backend and model-view registries & registry and model-view factory tests \\
Physical pair & differential-drive scenario and MJCF contact & physical docking/contact tests \\
Physical snake & physical reconfiguration and authored routes & full MuJoCo physical reconfiguration \\
\bottomrule
\end{longtable}

\section{Artifact Checklist}
\label{app:artifact}

For an archival release accompanying a future publication, include:

\begin{itemize}
  \item tagged commit, version, lock file, and built wheel/source distribution;
  \item explicit software LICENSE, asset license/NOTICE, and CITATION.cff;
  \item pinned TeX toolchain and rendered report PDF;
  \item exact command configurations and source-checkout scenarios;
  \item raw trace files with hashes;
  \item machine/OS/CPU/GPU disclosure in a separate benchmark manifest;
  \item videos synchronized with event sequence and simulated time;
  \item calibration data and hardware experimental protocol; and
  \item expected terminal metrics and tolerance bands.
\end{itemize}

